\chapter{Operating it}
\label{ch:operating}

\section*{What this chapter is for}

Observability, and testing a system that does not repeat. This chapter carries
a warning the others do not: \textbf{whether these questions get asked is
genuinely contested.}

Within one corpus, one analysis describes observability as implicit and
\reported{discussed but not heavily featured}{field-guide-system-design}, while
another describes it as \reported{non-negotiable}{field-guide-get-hired}. Two
files, one body of evidence, opposite conclusions.

\judgement{The usable resolution is not to guess which is right. It is to treat
this as material to raise unprompted rather than to wait to be asked \dash{}
which is the correct handling for a subject that is sometimes assessed and
always relevant.}

\section*{The questions}

\begin{bankq}{What would you instrument?}
\qasked
Asked as part of a design round, often near the end when the architecture is
settled and somebody asks how you would run it.

\qtested
Whether you know what is different about observing this kind of system.

\qstaff
Everything you already instrument, plus three things.

\textbf{The tokens.} Per request, input and output separately, because it is
the cost driver and because a change in average context length is the earliest
signal that something upstream changed.

\textbf{The decisions.} Which tools were called, with what arguments, in what
order, and how the turn ended. \judgement{This is the one that matters most and
the one usually missing. Without it a production failure is a paragraph of
prose and a shrug; with it, it is a trace you can replay, assert against, and
turn into an evaluation case.}

\textbf{The exact context.} What actually entered the model, not the identifiers
of what you intended to send. Chapter~\ref{ch:retrieval}'s diagnosis is
impossible without it, and it is the difference between debugging and guessing.

Two operational details worth naming, because both silently destroy the value
of the rest: sampling, which means most incidents have no trace at all; and
attribute truncation, which quietly removes exactly the long fields \dash{}
context, arguments \dash{} that you needed.

\qsenior
Latency, error rate, throughput, plus token counts. Correct and conventional,
and missing the decisions and the context, which are the two things specific to
this domain.

\qcontested
Whether the full context should be retained, given that it contains user data.
\judgement{A real tension with no general answer. Retention windows and
redaction at write time are the usual compromise, and the decision belongs to
the product rather than to engineering.}

\qdrill
For a system you know, ask whether you could reconstruct exactly what the model
saw on a request from last Tuesday. The answer is usually no.
\end{bankq}

\begin{bankq}{How do you regression-test something that does not repeat?}
\qasked
Sometimes asked directly. More often it arrives as \emph{how would you know if
a prompt change made it worse}, which is the same question.

\qtested
Whether you have a structure or an aspiration. It is Chapter~\ref{ch:evals}
from the operations side.

\qstaff
Two layers doing two jobs.

A deterministic layer over the behaviour \dash{} tool calls, order, arguments,
outcomes, absences \dash{} which needs no model, no credential and no network,
which is exactly why it can gate every change. This is where a regression is
caught.

A measured layer over quality, run on a schedule rather than per change,
compared against a recorded baseline. This is where a drift is noticed.

The baseline is the part that needs care: it is only meaningful against the
same cases and the same contract, so it records what it measured, and a
comparison across a changed case set or a changed rubric should be refused
rather than reported. \judgement{Silently comparing two numbers produced under
different conditions is worse than having no baseline, because it produces a
confident wrong answer.}

And a change to a prompt is a change to behaviour. It should go through review
with the same seriousness as a code change, and \judgement{it is worth having a
mechanical rule \dash{} a prompt edit requires a change to the specification or
the cases \dash{} because a behaviour change with no visible diff is otherwise
invisible to every process you have.}

\qsenior
A set of cases run before release, with output compared by similarity. It is
unstable, it cannot gate, and its failures cannot be attributed.

\qcontested
Whether a quality score should ever block a release.
\judgement{Chapter~\ref{ch:evals}'s position: the deterministic layer blocks,
the measured layer informs.}

\qdrill
For a model-driven feature you know, say what would go red if somebody changed
its instruction and made it worse. If nothing would, that is the gap.
\end{bankq}

\section*{What this chapter does not claim}

The contested status is stated at the top and is the most important thing in
the chapter. Two analyses within one corpus disagree, and no independent source
was found to settle it, so nothing here claims these questions will be asked.

The instrumentation list and the two-layer structure are the author's judgement,
consistent with the evaluation practice cited in Chapter~\ref{ch:evals}.
